% Independent review fragment; requires amsmath, amssymb, booktabs.
% This file documents the theorem instantiated by check_symmetric_source.py.
\section{Independent proof of the symmetric destructive guard}

The symmetric-guard 194-state variant is materially different from the
cached-controller 194-state machine: its primitive program has 3,646
rather than 3,649 instructions. The following proof checks its precise
source-level preservation. The physical refinement and additional state
fold are separate obligations, checked in the construction package.

A decrement instruction $\mathrm{dec}(r,z,p)$ jumps to $z$ when register
$r$ is zero; otherwise it decrements $r$ and jumps to $p$. Logical
registers $H,X$ have indices 5 and 6. The common source prefix is exactly
\[
\begin{array}{r|l@{\qquad}r|l}
0&\mathrm{dec}(H,4,1)&1&\mathrm{dec}(X,3,2)\\
2&\mathrm{jump}(0)&3&\mathrm{jump}(5)\\
4&\mathrm{dec}(X,10,5)&5&\mathrm{dec}(H,7,6)\\
6&\mathrm{jump}(5)&7&\mathrm{dec}(X,9,8)\\
8&\mathrm{jump}(5)&9&\mathrm{jump}(11)\\
10&\mathrm{halt}.&&
\end{array}
\]
Let $G(a,b)$ count source instructions until the first entry to control
10 or 11 from control zero and $(H,X)=(a,b)$. It does not count the halt
instruction itself. Every other register is arbitrary.

\paragraph{Guard lemma.}
The computation terminates at control 10 exactly when $a=b$, and otherwise
at control 11. In both cases $H=X=0$, and every other register is
unchanged. Its exact duration is
\[
G(a,b)=
\begin{cases}
3a+2,&a=b,\\
2a+b+4,&a>b,\\
3b+2,&a<b.
\end{cases}
\tag{SG1}
\]

\paragraph{Proof.}
If both registers are positive, instructions $0,1,2$ consume exactly
one unit from each and return to control zero in three steps. Induction
therefore reduces the pair to one with a zero component after
$3\min(a,b)$ steps. If both are zero, the two zero branches at 0 and 4
reach control 10.

The cleanup loop at control 5 has a separate exact invariant. With
$(H,X)=(u,v)$, each positive $H$ costs two steps ($5,6$) and reduces $H$
by one. Once $H=0$, each positive $X$ costs three steps ($5,7,8$) and
reduces $X$ by one. The final path $5,7,9$ exits at 11. Thus cleanup
always clears both registers in $2u+3v+3$ steps. This statement follows
by induction first on $u$ and then on $v$.

If $a>b$, after the paired decrements the path $0,1,3$ takes three
steps to cleanup with $(H,X)=(a-b-1,0)$. The total is
\[
3b+3+2(a-b-1)+3=2a+b+4.
\]
If $a<b$, the path $0,4$ takes two steps to cleanup with
$(H,X)=(0,b-a-1)$. The total is
\[
3a+2+3(b-a-1)+3=3b+2.
\]
All inspected instructions mention only $H$ and $X$, proving the full
register-frame statement. This establishes (SG1) for all natural inputs.
In particular the outcome and complete exit register vector are invariant
under exchanging the two guard inputs, although the duration can change.

\paragraph{Application to the changed Shift tail.}
Enter source control 3,638 with $(T,H,X)=(t,h,x)$ and arbitrary other
registers. The common prefix clears $T$, then transfers each unit of $H$
into both $T$ and $X$. It uses
\[
(2t+1)+(4h+1)=2t+4h+2
\]
steps and leaves $(T,H,X)=(h,0,k)$, where $k=h+x$.
The old tail restores $X$ to $H$, sets $X=1$, and jumps to control zero;
it uses $3k+3$ further steps and presents guard inputs $(k,1)$.
The new tail merely increments $H$ and jumps to zero; it uses two steps
and presents $(1,k)$. By the guard lemma, the two source computations
have exactly the same exit control and full register vector. They both
halt when $k=1$; otherwise they continue at source 11 with $H=X=0$.
No assumption $x=0$ is needed for this preservation lemma. The normal
cached-command interpretation uses the inherited $x=0$ command boundary,
where $k=1$ is precisely the intended test $h=1$.

The total old and new durations to first entry at 10 or 11 are
\[
D_{\rm old}=2t+4h+3k+5+G(k,1),\qquad
D_{\rm new}=2t+4h+4+G(1,k).
\]
Consequently
\[
D_{\rm old}-D_{\rm new}=
\begin{cases}
0,&k=0,\\
4,&k=1,\\
2k+4,&k\ge2.
\end{cases}
\tag{SG2}
\]
Counting the final halt instruction would add one to both equal-case
durations and leave (SG2) unchanged. This is a source-step comparison;
it is not a claim of a uniform literal-tape speed improvement.

\paragraph{Binding and global source preservation.}
The complete old and new maps agree at all controls $0,\ldots,3643$.
The old five-instruction tail at 3644--3648 is replaced by the two
instructions $\mathrm{inc}(H,3645)$ and $\mathrm{jump}(0)$.
The only entry from outside the Shift block 3638--3648 is the zero
branch from source 48 to source 3638; no outside instruction enters the
removed tail. These are exact finite properties checked by
\texttt{check\_symmetric\_source.py}, including every successor.
Thus equal configurations outside the changed block synchronize either
by the identical next instruction or by the finite positive block
simulation just proved. A block ending at 10 is followed by the same
halt instruction. Induction on synchronized segments and the bounded
trace argument therefore give halting equivalence of the two source
programs from source 25, for every initial register vector. All loop
inductions and finite trace constructions are available in PA.

\paragraph{Review evidence.}
The source checker accepts the exact literal maps and records their
hashes in \texttt{symmetric\_source\_certificate.json}. An independent
finite regression checks (SG1) on 4,225 guard pairs and (SG2), together
with full output equality, on 2,197 triples $(t,h,x)$. Its receipt is
\texttt{symmetric\_duration\_regression.json}. These finite checks catch
binding or arithmetic transcription errors; the universal justification
is the explicit induction above. No new proof-assistant elaboration is
claimed by this review.
